\begin{theorem}[有限深度族的代数闭包只记住 Fibonacci 模数的最小公倍数]\label{thm:conclusion-fibadic-finite-depth-family-algebra-closure}
设 \(\mathcal D\subset \NN_{\ge 1}\) 为有限非空集，并记
\[
L(\mathcal D):=\operatorname{lcm}\{F_e:\ e\in \mathcal D\}.
\]
定义含单位 \(^*\)-代数
\[
\mathcal A_{\mathcal D}
:=
\operatorname{Alg}^{*}\!\bigl(1,\ \bigcup_{e\in \mathcal D}\mathcal H_e\bigr)
\subset L^2(\widehat X_F,m_H),
\]
以及有限商代数
\[
A_N:=
\{f:\widehat X_F\to \CC:\ f=\widetilde f\circ \pi_N
\text{ for some }\widetilde f:\ZZ/(N\ZZ)\to\CC\}.
\]
则有严格恒等式
\[
\mathcal A_{\mathcal D}=A_{L(\mathcal D)}.
\]
特别地，
\[
\dim \mathcal A_{\mathcal D}=L(\mathcal D).
\]
\leanverified{paper\_conclusion\_fibadic\_finite\_depth\_family\_algebra\_closure}
\end{theorem}
\begin{proof}
每个 \(\mathcal H_e\) 都由经由 \(\pi_{F_e}\) 因子化的函数组成，故任意有限次线性组合、乘法与共轭之后所得函数必同时经由全部模数 \(F_e\) 的共同上界
\[
L(\mathcal D)=\operatorname{lcm}\{F_e:\ e\in \mathcal D\}
\]
因子化。于是
\[
\mathcal A_{\mathcal D}\subseteq A_{L(\mathcal D)}.
\]

反过来，对每个 \(e\in \mathcal D\)，\(\mathcal H_e\) 包含模 \(F_e\) 的本原角色
\[
\chi_{1/F_e}(x):=\exp\!\left(2\pi i\,\frac{\pi_{F_e}(x)}{F_e}\right).
\]
而 \(\QQ/\ZZ\) 中由 \(\{1/F_e\}_{e\in \mathcal D}\) 生成的子群恰为
\[
\frac{1}{L(\mathcal D)}\ZZ/\ZZ.
\]
故存在整数 \(u_e\) 使
\[
\sum_{e\in \mathcal D}\frac{u_e}{F_e}\equiv \frac{1}{L(\mathcal D)}\pmod{\ZZ},
\]
于是
\[
\prod_{e\in \mathcal D}\chi_{1/F_e}^{\,u_e}
=
\chi_{1/L(\mathcal D)}
\in
\mathcal A_{\mathcal D}.
\]
该角色的全部幂与共轭张成 \(\ZZ/(L(\mathcal D)\ZZ)\) 的全部角色，故其线性张成正是 \(A_{L(\mathcal D)}\)。于是
\[
A_{L(\mathcal D)}\subseteq \mathcal A_{\mathcal D}.
\]
两向合并即得结论。证毕。
\end{proof}

\begin{corollary}[真因子层生成的旧代数具有精确导体 \texorpdfstring{$Q_d$}{Q\_d}]\label{cor:conclusion-fibadic-old-algebra-exact-conductor-qd}
\leanverified{paper\_conclusion\_fibadic\_old\_algebra\_exact\_conductor\_qd}
对 \(d\ge 3\)，定义
\[
\mathcal D_d^{<}:=\{e:\ e\mid d,\ e<d\},
\qquad
\mathcal O_d:=\operatorname{Alg}^{*}\!\bigl(1,\ \bigcup_{e\in \mathcal D_d^{<}}\mathcal N_e\bigr).
\]
则有
\[
\mathcal O_d=A_{Q_d},
\qquad
Q_d=\operatorname{lcm}_{\substack{e\mid d\\ e<d}}F_e
=\frac{F_d}{\Pi_d(1)}.
\]
因而
\[
\dim \mathcal O_d=Q_d,
\qquad
\frac{\dim \mathcal H_d}{\dim \mathcal O_d}=\Pi_d(1).
\]
特别地，
\[
\mathcal O_d\subsetneq \mathcal H_d
\iff
\Pi_d(1)>1.
\]
\end{corollary}
\begin{proof}
由于
\[
\mathcal H_e=\bigoplus_{f\mid e}\mathcal N_f,
\]
且 \(\mathcal D_d^{<}\) 对整除向下封闭，故由 \(\{\mathcal N_e\}_{e\in \mathcal D_d^{<}}\) 生成的代数等同于由 \(\{\mathcal H_e\}_{e\in \mathcal D_d^{<}}\) 生成的代数。于是定理 \ref{thm:conclusion-fibadic-finite-depth-family-algebra-closure} 给出
\[
\mathcal O_d=A_{L_d},
\qquad
L_d:=\operatorname{lcm}_{\substack{e\mid d\\ e<d}}F_e.
\]
另一方面，定理 \ref{thm:conclusion-fibadic-conductor-budget-new-old-splitting} 已证明
\[
Q_d=\frac{F_d}{\Pi_d(1)}
=
\prod_{\substack{e\mid d\\ 1<e<d}}\Pi_e(1),
\]
而按同一定理的 \(p\)-进公式，
\[
v_p(Q_d)=\#\{k\ge 1:\ z(p^k)\mid d,\ z(p^k)<d\}
=v_p(L_d).
\]
故 \(L_d=Q_d\)。维数与严格包含判据随即成立。证毕。
\end{proof}

\begin{theorem}[精确深度 Ramanujan 包核在旧胞上严格失均值]\label{thm:conclusion-fibadic-depth-kernel-oldcell-invisibility}
对每个 \(d\ge 1\)，精确深度 Ramanujan 包核 \(K_d\) 满足
\[
P_df=f*K_d
\qquad
(f\in L^2(\widehat X_F,m_H)).
\]
若进一步 \(d\ge 3\) 且 \(\Pi_d(1)>1\)，记
\[
\mathbb E_{Q_d}:L^2(\widehat X_F,m_H)\to A_{Q_d}
\]
为到旧代数的正交投影，则
\[
\mathbb E_{Q_d}K_d=0.
\]
等价地，在有限商 \(\ZZ/(F_d\ZZ)\) 上，对每个 \(r\in \ZZ/(Q_d\ZZ)\) 都有
\[
\sum_{\substack{x\in \ZZ/(F_d\ZZ)\\ x\equiv r\;(\mathrm{mod}\ Q_d)}}K_d(x)=0.
\]
\end{theorem}
\begin{proof}
第一式已由定理 \ref{thm:conclusion-fibadic-ramanujan-kernel-convolution-primitive-idempotent} 给出。

现设 \(\Pi_d(1)>1\)。由推论 \ref{cor:conclusion-fibadic-old-algebra-exact-conductor-qd}，
\[
A_{Q_d}=\mathcal O_d
\]
正是全部真因子层生成的旧代数。若某个角色 \(\chi\) 满足 \(\delta(\chi)=d\) 且经由 \(\pi_{Q_d}\) 因子化，则其导体 \(N\) 必整除 \(Q_d\)。但 \(Q_d\) 的全部素幂只来自真因子深度，从而 \(\chi\) 的导体不可能仍具有精确深度 \(d\)，矛盾。故每个 \(\delta(\chi)=d\) 的角色在每个 \(Q_d\)-胞上的平均都为零。

把这些角色求和，便得到
\[
\mathbb E_{Q_d}K_d=0.
\]
离散分块求和式只是这一条件期望恒等式在有限商 \(\ZZ/(F_d\ZZ)\) 上的展开。证毕。
\end{proof}

\begin{theorem}[任意有限深度族的混合隐藏态最小受体律]\label{thm:conclusion-mixed-hidden-state-finite-depth-family-minimal-receiver}
固定可见屏幕 \(S\)，并记
\[
\beta:=\beta_n(\mathcal Q_m,A_{\neg S};\FF_2).
\]
对有限非空深度集 \(\mathcal D\subset \NN_{\ge 1}\)，记
\[
L(\mathcal D):=\operatorname{lcm}\{F_e:\ e\in \mathcal D\}.
\]
称有限阿贝尔群 \(G\) 为 \((S,\mathcal D)\)-忠实受体，若存在连续群同态
\[
\Theta:K_S^{\mathrm{col}}\times \widehat X_F\to G
\]
满足：
\begin{enumerate}
  \item \(\Theta|_{K_S^{\mathrm{col}}}\) 为单射；
  \item \(\Theta\) 在地址方向的核恰为 \(L(\mathcal D)\widehat{\ZZ}\)。
\end{enumerate}
则此类受体存在，当且仅当
\[
(\ZZ/2\ZZ)^\beta\oplus \ZZ/(L(\mathcal D)\ZZ)\hookrightarrow G.
\]
特别地，
\[
|G|_{\min}=2^\beta L(\mathcal D),
\]
并且达到该下界时，同构类型唯一：
\[
G\cong (\ZZ/2\ZZ)^\beta\oplus \ZZ/(L(\mathcal D)\ZZ).
\]
\end{theorem}
\begin{proof}
由定理 \ref{thm:conclusion-fibadic-finite-depth-family-algebra-closure}，\(\mathcal D\) 所生成的全部地址可观测代数恰为 \(A_{L(\mathcal D)}\)，故地址方向的最小忠实有限商正是
\[
\widehat X_F/L(\mathcal D)\widehat{\ZZ}\cong \ZZ/(L(\mathcal D)\ZZ).
\]
另一方面，碰撞方向依旧是
\[
K_S^{\mathrm{col}}\cong (\ZZ/2\ZZ)^\beta.
\]
因此 \((S,\mathcal D)\)-忠实观测等价于把有限截断混合态
\[
(\ZZ/2\ZZ)^\beta\times \ZZ/(L(\mathcal D)\ZZ)
\]
单射送入 \(G\)。现在应用定理 \ref{thm:conclusion-mixed-hidden-state-minimal-finite-receiver} 在导体
\[
N=L(\mathcal D)
\]
的情形，即得存在性判据、最小基数与唯一极小同构类型。证毕。
\end{proof}

\begin{corollary}[从全部真因子层到整层 \texorpdfstring{$d$}{d} 的不可约受体膨胀]\label{cor:conclusion-fibadic-old-to-full-depth-receiver-inflation}
对 \(d\ge 3\)，旧代数 \(\mathcal O_d\) 的最小混合受体基数为
\[
2^\beta Q_d,
\]
而整层 \(\mathcal H_d\) 的最小混合受体基数为
\[
2^\beta F_d.
\]
故从“全部真因子层的乘法闭包”提升到“整层 \(d\) 的全部可见观测”，所需最小有限阿贝尔受体恰发生乘法膨胀
\[
\frac{2^\beta F_d}{2^\beta Q_d}
=
\frac{F_d}{Q_d}
=
\Pi_d(1).
\]
特别地，当 \(\Pi_d(1)>1\) 时，此膨胀严格，且不能被任何真因子层组合与碰撞扭量重编码吸收。
\end{corollary}
\begin{proof}
对旧代数 \(\mathcal O_d=A_{Q_d}\)，由定理 \ref{thm:conclusion-mixed-hidden-state-finite-depth-family-minimal-receiver} 取 \(\mathcal D=\mathcal D_d^{<}\)，得最小受体基数为 \(2^\beta Q_d\)。对整层 \(\mathcal H_d=A_{F_d}\)，同理取 \(\mathcal D=\{d\}\)，得最小受体基数为 \(2^\beta F_d\)。两者相除并用推论 \ref{cor:conclusion-fibadic-old-algebra-exact-conductor-qd} 即得
\[
\frac{F_d}{Q_d}=\Pi_d(1).
\]
证毕。
\end{proof}

\begin{theorem}[Bernoulli 射流与全局 old/new 预算的精确互反]\label{thm:conclusion-fibadic-bernoulli-jet-oldnew-budget-reciprocity}
\leanverified{paper\_conclusion\_fibadic\_bernoulli\_jet\_oldnew\_budget\_reciprocity}
对 \(d>1\)，记
\[
S_2(d):=\sum_{e\mid d}\mu(d/e)F_e^2.
\]
则有精确恒等式
\[
2Q_d\,\Pi_d'(1)=a_dF_d,
\]
以及
\[
12Q_d\,\Pi_d''(1)
=
\bigl(S_2(d)+3a_d(a_d-2)\bigr)F_d.
\]
\end{theorem}
\begin{proof}
由推论 \ref{cor:conclusion-fibadic-old-algebra-exact-conductor-qd}，
\[
\Pi_d(1)=\frac{F_d}{Q_d}.
\]
再由推论 \ref{cor:conclusion-fibadic-depth-packet-slope-at-one}，
\[
2\,\Pi_d'(1)=a_d\,\Pi_d(1),
\]
代入即得
\[
2Q_d\,\Pi_d'(1)=a_dF_d.
\]

对二阶导数，令
\[
h_d(u):=\log \Pi_d(e^u).
\]
由定理 \ref{thm:conclusion-fibadic-depth-packet-bernoulli-mobius-jet}，
\[
h_d'(0)=\frac{a_d}{2},
\qquad
h_d''(0)=\frac{S_2(d)}{12}.
\]
另一方面，
\[
h_d'(0)=\frac{\Pi_d'(1)}{\Pi_d(1)},
\]
且
\[
h_d''(0)
=
\frac{\Pi_d''(1)}{\Pi_d(1)}
+\frac{\Pi_d'(1)}{\Pi_d(1)}
-\left(\frac{\Pi_d'(1)}{\Pi_d(1)}\right)^2.
\]
即
\[
h_d''(0)
=
\frac{\Pi_d''(1)}{\Pi_d(1)}
+\frac{\Pi_d'(1)}{\Pi_d(1)}
-\left(\frac{\Pi_d'(1)}{\Pi_d(1)}\right)^2.
\]
代入
\[
\frac{\Pi_d'(1)}{\Pi_d(1)}=\frac{a_d}{2}
\]
得
\[
\frac{\Pi_d''(1)}{\Pi_d(1)}
=
\frac{S_2(d)}{12}-\frac{a_d}{2}+\frac{a_d^2}{4}
=
\frac{S_2(d)+3a_d(a_d-2)}{12}.
\]
最后再乘以 \(\Pi_d(1)=F_d/Q_d\)，即得第二式。证毕。
\end{proof}

\begin{proposition}[新包 jet 的高阶增长：一阶近二次，二阶近三次]\label{prop:conclusion-fibadic-newpacket-jet-high-order-growth}
当 \(d\to\infty\) 时，
\[
S_2(d)=F_d^2+O\!\bigl(\tau(d)\varphi^d\bigr),
\]
从而
\[
\Pi_d'(1)=\left(\frac12+o(1)\right)\frac{F_d^2}{Q_d},
\qquad
\Pi_d''(1)=\left(\frac13+o(1)\right)\frac{F_d^3}{Q_d}.
\]
若再假设
\[
\log Q_d=o(d),
\]
则进一步有
\[
\Pi_d'(1)=\frac12\,F_d^{\,2-o(1)},
\qquad
\Pi_d''(1)=\frac13\,F_d^{\,3-o(1)}.
\]
\end{proposition}
\begin{proof}
由定义
\[
S_2(d)=\sum_{e\mid d}\mu(d/e)F_e^2
=
F_d^2+\sum_{\substack{e\mid d\\ e<d}}\mu(d/e)F_e^2.
\]
任意真因子 \(e<d\) 满足 \(e\le d/2\)，故
\[
\left|S_2(d)-F_d^2\right|
\le
\sum_{\substack{e\mid d\\ e<d}}F_e^2
=
O\!\bigl(\tau(d)\varphi^d\bigr).
\]
另一方面，文稿已给出 \(a_d/F_d\to 1\)。把这两条估计代入定理 \ref{thm:conclusion-fibadic-bernoulli-jet-oldnew-budget-reciprocity}，便得到
\[
\Pi_d'(1)=\frac{a_dF_d}{2Q_d}
=
\left(\frac12+o(1)\right)\frac{F_d^2}{Q_d},
\]
以及
\[
\Pi_d''(1)
=
\frac{S_2(d)+3a_d(a_d-2)}{12}\cdot \frac{F_d}{Q_d}
=
\left(\frac13+o(1)\right)\frac{F_d^3}{Q_d}.
\]
若再用 \(\log Q_d=o(d)\) 且 \(\log F_d\asymp d\)，则
\[
Q_d=F_d^{o(1)},
\]
从而上式化为
\[
\Pi_d'(1)=\frac12\,F_d^{\,2-o(1)},
\qquad
\Pi_d''(1)=\frac13\,F_d^{\,3-o(1)}.
\]
证毕。
\end{proof}

\begin{conjecture}[旧胞二次能量均分]\label{conj:conclusion-fibadic-oldcell-quadratic-energy-equipartition}
设 \(d\to\infty\) 且 \(\Pi_d(1)>1\)。在有限商 \(\ZZ/(F_d\ZZ)\) 中，以模 \(Q_d\) 的剩余类分解为旧胞
\[
C_r:=\{x\in \ZZ/(F_d\ZZ): x\equiv r\pmod{Q_d}\}.
\]
则精确深度核 \(K_d\) 应满足二次能量均分律
\[
\max_{r\in \ZZ/(Q_d\ZZ)}
\left|
\sum_{x\in C_r}|K_d(x)|^2
-
\frac{F_d}{Q_d}\,a_d
\right|
=
o\!\left(\frac{F_d}{Q_d}\,a_d\right).
\]
\end{conjecture}

\noindent
这条链把 old/new 分裂从线性层推进到代数与硬件层：真因子层的全部乘法闭包恰止于导体 \(Q_d\) 的旧代数，精确深度包核在每个旧胞内严格失均值，而从旧代数跃迁到整层 \(\mathcal H_d\) 的最小有限受体又恰按 \(\Pi_d(1)\) 发生乘法膨胀；与此同时，这一全局膨胀又被 \(\Pi_d\) 在 \(X=1\) 处的一阶、二阶 jet 精确编码。
